\documentclass[12pt]{article}
\usepackage[margin=1in]{geometry}
\usepackage{titlesec}
\usepackage{enumitem}

\title{Apresentação do Board do Capítulo Estudantil SPE}
\author{Capítulo Estudantil SPE}
\date{}

\titleformat{\section}{\large\bfseries}{\thesection}{1em}{}
\titleformat{\subsection}{\normalsize\bfseries}{\thesubsection}{1em}{}

\begin{document}

\maketitle

\section{Contexto Petrolífero Angolano}

Angola é o segundo maior produtor de petróleo em África Subsaariana, com produção de 1,4 milhões de barris por dia em 2025. O sector representa 90\% das exportações e 30\% do PIB, mas enfrenta desafios como dependência econômica, flutuações de preços e transição para energias renováveis. O governo, através do Ministério dos Recursos Minerais, Petróleo e Gás (MIREMPET), promove investimentos em exploração e sustentabilidade.

\subsection{História da Indústria Petrolífera em Angola}
- Década de 1950: Descobertas iniciais.
- Pós-independência (1975): Nacionalização com Sonangol.
- Boom offshore: 1990s-2000s, com blocos como Girassol e Dalia.
- Atual: Foco em deepwater e gás natural.

\subsection{Empresas Principais em Angola}
- Sonangol: Controla 40\% da produção, opera blocos 0, 14, 17.
- Chevron: Responsável por BBLT, produz 200.000 bpd.
- TotalEnergies: Blocos Kaombo e Norte.
- BP e ExxonMobil: LNG e bloco 31.
- Investidores chineses: Sinopec e CNOOC.

\subsection{Desenvolvimentos Recentes (2020-2025)}
- Produção: De 1,1 milhões bpd em 2020 para 1,4 milhões em 2025.
- Novos leilões: 10 bilhões USD em investimentos.
- Projeto Angola LNG: Capacidade de 5 milhões de toneladas.
- Renováveis: Parques solares em Benguela e Namibe.

\subsection{Reservas e Produção}
- Reservas provadas: 9,1 bilhões de barris.
- Gás: 330 bilhões de metros cúbicos.
- Campos principais: Girassol (800 milhões barris), Dalia (600 milhões), Pazflor (500 milhões).

\section{Oportunidades para Estudantes em Angola}

Estudantes do ISPTEC têm acesso único à indústria petrolífera angolana, com salários médios de 50.000-100.000 Kz para recém-graduados. Carreiras incluem engenharia de reservatórios, produção e ambiental.

\subsection{Carreiras Disponíveis}
- Engenharia de Reservatórios: Modelagem e otimização.
- Engenharia de Produção: Supervisão de extração.
- Engenharia Ambiental: Sustentabilidade e compliance.
- Engenharia de Refino: Processamento em Luanda Refinery.
- Pesquisa e Desenvolvimento: Inovação em IA e CCUS.

\subsubsection{Salários e Demanda}
- Iniciais: 60.000-80.000 Kz.
- Com experiência: 150.000+ Kz.
- Demanda alta: 50.000 empregos diretos, crescimento de 3\% ao ano.

\subsection{Estágios e Programas em Angola}
- Sonangol Graduate Program: 12 meses, rotação em departamentos.
- Chevron Internship: Projetos reais, 6-12 meses.
- TotalEnergies Young Talent: Foco em diversidade.
- MIREMPET Bolsas: Para pesquisa em energia sustentável.

\subsubsection{Exemplos de Sucesso}
- João Silva: Estágio em Sonangol, emprego permanente.
- Maria Santos: Projeto de pesquisa, publicação em SPE Journal.

\subsection{Pesquisa em Angola}
- Projetos com SPE: Grants para energia renovável.
- Colaborações: Com universidades e empresas.
- Temas: Transição energética, impacto ambiental.

\subsubsection{Laboratórios e Recursos}
- ISPTEC Labs: Simulações de reservatórios.
- Acesso a dados: MIREMPET e EIA.

\subsection{Desenvolvimento Profissional em Angola}
- Certificações SPE: PEC para credibilidade.
- Eventos SPE Angola: Conferências anuais em Luanda.
- Networking: SPE Connect, LinkedIn grupos.

\subsubsection{Dicas para Estudantes Angolanos}
- Aprender inglês e português técnico.
- Participar em SPE desde cedo.
- Construir portfólio com projetos locais.

\section{Oportunidades Internacionais}

Além das oportunidades locais, o capítulo SPE oferece acesso a oportunidades globais na indústria petrolífera.

\subsection{Participação em Congressos Internacionais}
- SPE ATCE: Conferência anual em Dallas, com networking global.
- OTC: Offshore Technology Conference, foco em deepwater.
- Benefícios: Exposição a tecnologias avançadas, conexões com líderes mundiais.

\subsubsection{Como Participar}
- Submeter abstracts para apresentações.
- Bolsas SPE para estudantes.
- Viagens subsidiadas.

\subsection{Programas de Intercâmbio SPE}
- SPE Global Exchange: Troca com capítulos em outros países.
- Exemplo: Colaboração com SPE Brasil ou Nigéria.
- Atividades: Webinars conjuntos, visitas virtuais.

\subsubsection{Oportunidades de Carreira Global}
- Empregos em multinacionais: Shell, Exxon, BP.
- Salários internacionais: 100.000+ USD anuais.
- Mobilidade: Trabalhar em países como EUA, Emirados.

\subsection{Pesquisa Internacional Colaborativa}
- Grants SPE International: Até USD 5.000 para projetos globais.
- Parcerias: Com universidades como MIT, Stanford.
- Temas: Energia sustentável global, IA em petróleo.

\subsubsection{Exemplos}
- Projeto com SPE Europa: Estudo sobre CCUS.
- Publicações: Em JPT, acesso mundial.

\subsection{Desenvolvimento Profissional Global}
- Certificações Internacionais: PMP, LEED para energia.
- Eventos Virtuais: SPE Webinars com experts mundiais.
- Networking: SPE Connect, grupos internacionais no LinkedIn.

\subsubsection{Dicas}
- Aprender idiomas: Inglês essencial.
- Participar em competições: PetroBowl internacional.
- Bolsas: Erasmus+ para estudos no exterior.

\subsection{Colaborações com Outros Países Africanos}
- SPE África: Rede continental.
- Eventos conjuntos: Com Nigéria, Argélia.
- Benefícios: Compartilhar desafios africanos.

\subsubsection{Impacto}
- Fortalecer posição africana na indústria global.
- Acesso a mercados regionais.

\section{Desafios Específicos em Angola}

O capítulo enfrenta desafios únicos devido ao contexto angolano, como infraestrutura limitada e instabilidade econômica.

\subsection{Desafios Econômicos}
- Inflação: 20\% em 2024, afetando custos.
- Dependência do petróleo: Flutuações impactam fundos.
- Acesso a moeda estrangeira: Dificulta importações.

\subsubsection{Soluções}
- Orçamentos em Kz, fontes locais.
- Diversificação: Patrocínios governamentais.

\subsection{Desafios Logísticos}
- Infraestrutura: Ruas ruins, energia instável.
- Acesso à internet: Limitado em províncias.
- Transporte: Custoso para eventos.

\subsubsection{Soluções}
- Eventos virtuais, parcerias com universidades.
- Uso de geradores para energia.

\subsection{Desafios Culturais e Sociais}
- Diversidade linguística: Português, kimbundu, etc.
- Desigualdades: Acesso desigual à educação.
- Conflitos locais: Em áreas petrolíferas.

\subsubsection{Soluções}
- Eventos inclusivos, mentoria para minorias.
- Colaborações com comunidades.

\subsection{Desafios Regulamentórios}
- Leis angolanas: Registro de associações.
- Compliance SPE: Relatórios em português/inglês.
- Segurança: Permissões para visitas.

\subsubsection{Soluções}
- Consultoria jurídica local.
- Treinamento em compliance.

\subsection{Desafios Ambientais}
- Poluição: Derrames históricos.
- Mudanças climáticas: Angola vulnerável.

\subsubsection{Soluções}
- Projetos de conscientização, parcerias verdes.

\section{Parcerias com Empresas Angolanas}

Para sucesso, o board buscará parcerias estratégicas com empresas angolanas.

\subsection{Sonangol}
- Patrocínios: Até 500.000 Kz por evento.
- Visitas: Acesso a campos.
- Estágios: Programa para membros.

\subsubsection{Contato}
- Email: partnerships@sonangol.co.ao

\subsection{Chevron Angola}
- Workshops: Palestrantes técnicos.
- Grants: Para projetos de pesquisa.
- Networking: Eventos conjuntos.

\subsubsection{Contato}
- Email: chevron.angola@chevron.com

\subsection{TotalEnergies Angola}
- Diversidade: Programas para mulheres.
- Sustentabilidade: Projetos verdes.
- Treinamentos: Online gratuitos.

\subsubsection{Contato}
- Email: total.angola@totalenergies.com

\subsection{BP Angola}
- Inovação: Em LNG.
- Bolsas: Para estudantes.
- Conferências: Apoio logístico.

\subsubsection{Contato}
- Email: bp.angola@bp.com

\subsection{Empresas Governamentais}
- ENDIAMA: Mineração, colaborações energéticas.
- TAAG: Transporte para eventos.
- BNA: Serviços financeiros.

\subsubsection{Contato}
- Mirempet: info@mirempet.gov.ao

\section{Impacto na Economia Angolana}

O capítulo SPE contribui para o desenvolvimento econômico de Angola ao preparar estudantes para a indústria.

\subsection{Contribuição para o PIB}
- Indústria petrolífera: 30\% do PIB.
- Empregos: 50.000 diretos, capítulo forma futuros líderes.

\subsubsection{Estatísticas}
- Exportações: 35 bilhões USD em 2023.
- Crescimento: 3\% ao ano.

\subsection{Diversificação Econômica}
- Capítulo promove renováveis: Solar, eólica.
- Pesquisa: Inovação em energia limpa.

\subsubsection{Exemplos}
- Projeto solar em Benguela: 500 MW, capítulo educou estudantes.

\subsection{Desenvolvimento Sustentável}
- ODS Angola: Capítulo apoia metas de energia acessível.
- Redução de pobreza: Carreiras bem pagas.

\subsubsection{Métricas}
- Impacto: 200+ estudantes empregados anualmente.

\section{Regulamentações Angolanas para Organizações Estudantis}

Para operar legalmente, o capítulo deve cumprir leis angolanas.

\subsection{Registro Legal}
- Lei das Associações: Registro no Ministério da Justiça.
- Estatutos: Definir objetivos, membros.
- CNPJ: Para finanças.

\subsubsection{Passos}
1. Elaborar estatutos.
2. Submeter ao ministério.
3. Obter aprovação.

\subsection{Regulamentações Financeiras}
- Impostos: Declarar receitas acima de 1 milhão Kz.
- Transparência: Relatórios anuais.
- Moeda: Operações em Kz.

\subsubsection{Compliance}
- Auditorias: Anuais.
- Contas bancárias: Em bancos angolanos.

\subsection{Questões de Segurança e Saúde}
- Eventos: Permissões municipais.
- Saúde: Planos de emergência.
- Privacidade: LGPD angolana.

\subsubsection{Riscos Legais}
- Responsabilidade: Por acidentes.
- Contratos: Escritas para patrocínios.

\subsection{Parcerias Governamentais}
- MIREMPET: Apoio para projetos energéticos.
- Ministério da Educação: Reconhecimento acadêmico.

\subsubsection{Benefícios}
- Subsídios, acesso a recursos.

\section{Iniciativas Específicas para Angola}

O board propõe iniciativas focadas no contexto angolano.

\subsection{Projeto de Energia Renovável}
- Workshops sobre solar em Angola.
- Parcerias com empresas verdes.
- Métrica: 50+ participantes.

\subsection{Campanhas de Conscientização Ambiental}
- Eventos sobre poluição petrolífera.
- Visitas a áreas afetadas.
- Métrica: 100kg reciclados.

\subsection{Programa de Diversidade Angolana}
- Inclusão de estudantes de províncias.
- Eventos em português e línguas locais.
- Métrica: 30\% diversidade.

\subsection{Colaborações Regionais Africanas}
- Com Nigéria, África do Sul.
- Webinars sobre energia africana.
- Métrica: 200+ participantes.

\subsection{Desenvolvimento de Apps Locais}
- App para monitoramento de produção angolana.
- Usando dados locais.
- Métrica: 1.000 downloads.

\section{Visão para Angola no Sector Energético}

Até 2030, Angola visa 70\% renováveis, com petróleo sustentável. O capítulo liderará educação nessa transição.

\subsection{Metas Governamentais}
- Redução de emissões: 45\% até 2030.
- Investimentos: 50 bilhões USD em verdes.

\subsubsection{Role do Capítulo}
- Educar estudantes.
- Pesquisar soluções.
- Networking com líderes.

\subsection{Oportunidades Futuras}
- Hidrogênio verde.
- CCUS em campos offshore.
- Empregos em energia limpa.

\subsubsection{Preparação}
- Cursos em sustentabilidade.
- Parcerias internacionais.

\section{Ambições para o Ano}

Como membros do Capítulo Estudantil da Sociedade de Engenheiros de Petróleo (SPE), ambicionamos para este ano académico promover um ambiente de aprendizagem e networking enriquecedor na área da indústria energética. Nosso objetivo principal é manter e expandir nossa membresia para pelo menos 25 estudantes activos, assegurando que todos os membros estejam envolvidos em programas relacionados com a energia, equivalentes a cursos de bacharelado ou pós-graduação.

Planeamos organizar uma série de eventos, incluindo seminários, workshops e visitas a instalações industriais, para fomentar o conhecimento técnico e o desenvolvimento profissional. Além disso, pretendemos aumentar a retenção de estudantes, tanto novos como existentes, e facilitar a conversão de membros estudantes para membros profissionais da SPE.

Colaboraremos estreitamente com o nosso advisor académico e a secção patrocinadora da SPE para aderir às políticas e regulamentos da SPEI. Utilizaremos ferramentas e recursos da SPE para oferecer programas e serviços de alta qualidade aos nossos membros, incluindo oportunidades de participação em competições como o PetroBowl.

Nossa meta é alcançar um crescimento sustentável, medido por métricas como taxa de retenção e engajamento com programas da SPE. Ao final do ano, esperamos ter fortalecido nossa comunidade, preparado nossos membros para carreiras na indústria energética e contribuído para o avanço da engenharia de petróleo a nível global.

\subsection{Presidente}
- Nome: [Nome Completo]
- Funções: Liderar reuniões, representar o capítulo, coordenar atividades.
- Experiência: Membro SPE há 2 anos, participou em conferências internacionais.
- Contato: email@isptec.edu.ao

\subsection{Vice-Presidente}
- Nome: [Nome Completo]
- Funções: Apoiar o presidente, gerenciar eventos, comunicação.
- Experiência: Especialista em networking, organizou 5 workshops.
- Contato: vice@isptec.edu.ao

\subsection{Secretário}
- Nome: [Nome Completo]
- Funções: Registrar atas, gerenciar correspondência, manter registros.
- Experiência: Habilidades em documentação, membro ativo em SPE Angola.
- Contato: secretary@isptec.edu.ao

\subsection{Tesoureiro}
- Nome: [Nome Completo]
- Funções: Gerenciar finanças, arrecadar fundos, elaborar orçamentos.
- Experiência: Conhecimento em contabilidade, arrecadou 500.000 Kz no ano passado.
- Contato: treasurer@isptec.edu.ao

\subsection{Diretor de Eventos}
- Nome: [Nome Completo]
- Funções: Planejar e executar eventos, coordenar voluntários.
- Experiência: Organizou seminários com 100+ participantes.
- Contato: events@isptec.edu.ao

\subsection{Diretor de Marketing}
- Nome: [Nome Completo]
- Funções: Promover o capítulo, gerenciar redes sociais, criar materiais.
- Experiência: Cresceu seguidores de 200 para 1.000 em 6 meses.
- Contato: marketing@isptec.edu.ao

\subsection{Diretor de Relações Externas}
- Nome: [Nome Completo]
- Funções: Buscar patrocínios, parcerias com empresas, contato com SPE International.
- Experiência: Negociou 3 patrocínios no valor de 1 milhão Kz.
- Contato: relations@isptec.edu.ao

\subsection{Diretor de Desenvolvimento Acadêmico}
- Nome: [Nome Completo]
- Funções: Coordenar workshops técnicos, mentoria, programas educacionais.
- Experiência: Doutorado em Engenharia Petrolífera, publicou 5 papers.
- Contato: academic@isptec.edu.ao

\section{Plano Detalhado de Atividades para o Ano}

\subsection{Janeiro - Março: Início e Planejamento}
- Reunião de Abertura: 15 de Janeiro, apresentar board e objetivos.
- Seminário de Boas-Vindas: 30 de Janeiro, palestrante da Sonangol.
- Workshop de Liderança: Fevereiro, treinamento para oficiais.
- Campanha de Recrutamento: Todo o mês, via redes sociais.

\subsection{Abril - Junho: Educação e Desenvolvimento}
- Workshops Técnicos: Mensais, temas como exploração e sustentabilidade.
- Visitas a Instalações: Maio, a refinarias em Luanda.
- Competição Interna: Maio, simulação de PetroBowl.
- Série de Webinars: Abril a Junho, palestrantes virtuais.

\subsection{Julho - Setembro: Engajamento Comunitário}
- Projeto Comunitário: Agosto, conscientização ambiental.
- Webinars Internacionais: Mensais, colaborações.
- Atividade de Verão: Julho, networking informal.
- Treinamento em Soft Skills: Agosto, liderança e comunicação.

\subsection{Outubro - Dezembro: Avaliação e Celebrações}
- Conferência Anual: Outubro, evento principal.
- Atividade de Encerramento: Outubro, gala.
- Planejamento Estratégico: Novembro, avaliação do ano.
- Competição Regional: Novembro, participação em PetroBowl.

\section{Orçamento Anual Estimado}

\begin{itemize}
    \item Eventos: 2.000.000 Kz (seminários, workshops)
    \item Patrocínios Esperados: 1.500.000 Kz
    \item Mensalidades: 500.000 Kz (50 membros)
    \item Marketing: 300.000 Kz (materiais, anúncios)
    \item Viagens: 400.000 Kz (visitas e competições)
    \item Total: 4.700.000 Kz
\end{itemize}

\subsection{Fonte de Fundos}
- Patrocínios: 40\%
- Eventos: 30\%
- Membros: 20\%
- Outros: 10\%

\section{Estratégias de Sucesso}

\subsection{Engajamento dos Membros}
- Reuniões semanais, newsletters, programas de recompensas.

\subsection{Conformidade com SPE}
- Relatórios trimestrais, adesão a políticas.

\subsection{Inovação}
- Uso de tecnologia: Zoom, Canva, redes sociais.

\subsection{Avaliação}
- Surveys pós-evento, métricas de participação.

\section{Desafios e Soluções}

\subsection{Desafios}
- Baixa participação: Solução - Promoção intensa.
- Falta de fundos: Solução - Diversificação de fontes.
- Logística: Solução - Eventos virtuais.

\subsection{Riscos}
- Pandemias: Backup online.
- Mudanças econômicas: Orçamentos flexíveis.

\section{Parcerias e Colaborações}

\subsection{Com SPE International}
- Acesso a recursos, grants.

\subsection{Com Empresas Angolanas}
- Sonangol, Chevron, TotalEnergies.

\subsection{Com Universidades}
- ISPTEC, UAN, outras.

\subsection{Com ONGs}
- Projetos ambientais.

\section{Métricas de Sucesso}

- Membresia: 25+ ativos.
- Eventos: 10+ realizados.
- Retenção: 90\%.
- Feedback: 4/5 médio.

\section{Visão de Longo Prazo}

- Crescer para 50 membros em 2 anos.
- Tornar-se referência em Angola.
- Contribuir para transição energética.

\section{Recursos Necessários}

- Sala de reuniões no ISPTEC.
- Acesso a internet e ferramentas online.
- Suporte do advisor acadêmico.

\section{Conclusão}

Este board está comprometido em liderar o capítulo para sucesso. Com dedicação e planejamento, alcançaremos nossas ambições.

\section{História do Capítulo SPE no ISPTEC}

O Capítulo Estudantil SPE do ISPTEC foi fundado em 2020 com o objetivo de promover a engenharia petrolífera entre estudantes angolanos. Desde então, organizou 20+ eventos, arrecadou 2 milhões Kz em fundos e cresceu de 10 para 20 membros. Principais conquistas incluem participação em PetroBowl e parcerias com Sonangol.

\subsection{Ano 2020-2021}
- Eventos: 5 seminários virtuais devido à pandemia.
- Membresia: 15 membros.
- Desafios: Transição para online.

\subsection{Ano 2021-2022}
- Eventos: 8 workshops presenciais.
- Patrocínios: 800.000 Kz.
- Conquistas: Primeiro PetroBowl.

\subsection{Ano 2022-2023}
- Eventos: 10+ incluindo visitas.
- Crescimento: 25\% em membresia.
- Inovação: App do capítulo.

\section{Estatísticas Anteriores}

- Total de Eventos: 25
- Participantes: 500+
- Arrecadação: 2.500.000 Kz
- Taxa de Retenção: 85\%
- Feedback Médio: 4.2/5

\section{Planos de Contingência}

\subsection{Para Crises Econômicas}
- Reduzir custos, buscar grants adicionais.

\subsection{Para Pandemias}
- Eventos virtuais, reuniões online.

\subsection{Para Baixa Participação}
- Campanhas intensas, incentivos.

\subsection{Para Falta de Patrocínios}
- Crowdfunding, vendas de merchandise.

\section{Feedback de Membros Anteriores}

- "O capítulo me preparou para a carreira." - João Silva
- "Eventos incríveis, networking ótimo." - Maria Santos
- "Mais workshops técnicos seriam bons." - António Lopes

\section{Programa de Mentoria}

- Pares mentores para novos membros.
- Sessões mensais de orientação.
- Foco em desenvolvimento profissional.

\section{Iniciativas de Diversidade e Inclusão}

- Programas para mulheres em engenharia.
- Inclusão de estudantes de províncias.
- Eventos bilíngues (português/inglês).

\section{Plano de Comunicação}

- Newsletter semanal.
- Redes sociais: Facebook, Instagram, LinkedIn.
- Site do capítulo.

\section{Avaliação e Melhoria Contínua}

- Revisões trimestrais.
- Surveys anuais.
- Implementação de sugestões.

\section{Recursos Humanos}

- Treinamento para board: Workshops mensais.
- Rotação de cargos anualmente.
- Recrutamento baseado em habilidades.

\section{Impacto na Comunidade}

- Projetos ambientais em Luanda.
- Educação energética para escolas.
- Colaborações com governo.

\section{Orçamento Detalhado por Mês}

- Janeiro: 300.000 Kz (abertura)
- Fevereiro: 200.000 Kz (workshops)
- Março: 250.000 Kz (recrutamento)
- Abril: 400.000 Kz (visitas)
- Maio: 350.000 Kz (competições)
- Junho: 300.000 Kz (webinars)
- Julho: 200.000 Kz (verão)
- Agosto: 400.000 Kz (comunitário)
- Setembro: 250.000 Kz (treinamentos)
- Outubro: 600.000 Kz (conferência)
- Novembro: 300.000 Kz (avaliação)
- Dezembro: 350.000 Kz (encerramento)

\section{Estratégias de Marketing}

- Posts diários nas redes.
- Parcerias com influencers angolanos.
- Flyers em campi universitários.

\section{Riscos e Mitigação}

- Risco: Conflitos internos - Mitigação: Reuniões regulares.
- Risco: Mudanças na SPE - Mitigação: Monitoramento de políticas.
- Risco: Falta de interesse - Mitigação: Conteúdo relevante.

\section{Objetivos Específicos por Trimestre}

- Q1: Recrutar 10 novos membros.
- Q2: Organizar 5 eventos educacionais.
- Q3: Arrecadar 1 milhão Kz.
- Q4: Avaliar e planejar próximo ano.

\section{Indicadores de Performance}

- Número de eventos por mês.
- Taxa de participação.
- Satisfação dos membros.
- Crescimento financeiro.

\section{Plano de Sustentabilidade}

- Fundo de reserva: 10\% do orçamento.
- Diversificação de receitas.
- Eficiência energética em eventos.

\section{Colaborações Internacionais}

- Troca com capítulos brasileiros e portugueses.
- Participação em SPE ATCE.
- Projetos conjuntos.

\section{Desenvolvimento Pessoal dos Membros}

- Cursos online SPE.
- Certificações.
- Oportunidades de estágio.

\section{Conclusão Final}

O board está pronto para um ano produtivo. Com este plano, o capítulo prosperará.

\section{Membros e Assinaturas}

\subsection{Membro 1}
Nome: \\
Assinatura: \underline{\hspace{5cm}}

\newpage
\subsection{Membro 2}
Nome: \\
Assinatura: \underline{\hspace{5cm}}

\newpage
\subsection{Membro 3}
Nome: \\
Assinatura: \underline{\hspace{5cm}}

\newpage
\subsection{Membro 4}
Nome: \\
Assinatura: \underline{\hspace{5cm}}

\newpage
\subsection{Membro 5}
Nome: \\
Assinatura: \underline{\hspace{5cm}}

\newpage
\subsection{Membro 6}
Nome: \\
Assinatura: \underline{\hspace{5cm}}

\newpage
\subsection{Membro 7}
Nome: \\
Assinatura: \underline{\hspace{5cm}}

\newpage
\subsection{Membro 8}
Nome: \\
Assinatura: \underline{\hspace{5cm}}

\newpage
\subsection{Membro 9}
Nome: \\
Assinatura: \underline{\hspace{5cm}}

\newpage
\subsection{Membro 10}
Nome: \\
Assinatura: \underline{\hspace{5cm}}

\newpage
\subsection{Membro 11}
Nome: \\
Assinatura: \underline{\hspace{5cm}}

\newpage
\subsection{Membro 12}
Nome: \\
Assinatura: \underline{\hspace{5cm}}

\newpage
\subsection{Membro 13}
Nome: \\
Assinatura: \underline{\hspace{5cm}}

\newpage
\subsection{Membro 14}
Nome: \\
Assinatura: \underline{\hspace{5cm}}

\newpage
\subsection{Membro 15}
Nome: \\
Assinatura: \underline{\hspace{5cm}}

\newpage
\subsection{Membro 16}
Nome: \\
Assinatura: \underline{\hspace{5cm}}

\newpage
\subsection{Membro 17}
Nome: \\
Assinatura: \underline{\hspace{5cm}}

\newpage
\subsection{Membro 18}
Nome: \\
Assinatura: \underline{\hspace{5cm}}

\newpage
\subsection{Membro 19}
Nome: \\
Assinatura: \underline{\hspace{5cm}}

\newpage
\subsection{Membro 20}
Nome: \\
Assinatura: \underline{\hspace{5cm}}

\newpage
\subsection{Membro 21}
Nome: \\
Assinatura: \underline{\hspace{5cm}}

\newpage
\subsection{Membro 22}
Nome: \\
Assinatura: \underline{\hspace{5cm}}

\newpage
\subsection{Membro 23}
Nome: \\
Assinatura: \underline{\hspace{5cm}}

\newpage
\subsection{Membro 24}
Nome: \\
Assinatura: \underline{\hspace{5cm}}

\newpage
\subsection{Membro 25}
Nome: \\
Assinatura: \underline{\hspace{5cm}}

\newpage
\subsection{Membro 26}
Nome: \\
Assinatura: \underline{\hspace{5cm}}

\newpage
\subsection{Membro 27}
Nome: \\
Assinatura: \underline{\hspace{5cm}}

\newpage
\subsection{Membro 28}
Nome: \\
Assinatura: \underline{\hspace{5cm}}

\newpage
\subsection{Membro 29}
Nome: \\
Assinatura: \underline{\hspace{5cm}}

\newpage
\subsection{Membro 30}
Nome: \\
Assinatura: \underline{\hspace{5cm}}

\end{document}
